In this work we have developed a framework for calibrating tomographic redshift distributions for LSST-like weak-lensing and large-scale-structure analyses. This represents the latest step in a programme of augmentation-based redshift calibration: \citet{2024ApJ...967L...6M} first demonstrated the effectiveness of supplementing spectroscopic training sets with simulated galaxies as a proof of concept for LSST; Paper~I extended this into a SOM-based framework that enables flexible identification of under-represented regions of colour space and validated the approach under realistic LSST Year~1 and Year~10 scenarios against science requirements; and the present paper carries the augmentation framework from galaxy-level photometric-redshift estimation to population-level redshift-distribution inference. Our construction combines an importance-weighted Bayesian mixture formalism with a SOM-based discretisation of photometric feature space, enabling the redshift distribution in each tomographic bin to be written as a weighted sum of cluster-level conditional densities. Within this common representation, we compared direct spectroscopic summarisation (\texttt{DIR}), posterior stacking (\texttt{Stack}), and a geometric \texttt{Hybrid} estimator, and assessed their behaviour across multiple calibration configurations for both LSST-Y1 and LSST-Y10 survey scenarios.

Several conclusions emerge from these tests. First, lens-galaxy calibration is comparatively robust: for this brighter and narrower-redshift population, all estimators perform well, and the residual sensitivity to configuration choices is modest except in the highest-redshift bins. Secondly, source-galaxy calibration is substantially more demanding, owing to stronger colour--redshift degeneracies, weaker high-redshift support, and greater sensitivity to the representativeness of the calibration sample. In this regime, the \texttt{Hybrid} estimator provides the best overall compromise. It suppresses the excessive broadening often seen in \texttt{DIR}, while avoiding the over-regularisation and tail suppression that can arise in \texttt{Stack}, and therefore yields the most stable and least biased reconstructions across the physically motivated configurations. The incorporation of SOM-based augmentation is particularly beneficial for the source samples: it improves tomographic assignments, enlarges the supported region of feature space, and reduces the severity of high-redshift tail truncation. Under the best-performing augmentation-enabled configurations, the expectation-value biases can be reduced to the sub-percent level for lens samples and for most source bins in both the Y1 and Y10 scenarios, although the highest-redshift source bins remain the most challenging regime.

A central contribution of this paper is the propagation of calibration uncertainty beyond any single fixed experiment. By sampling across an ensemble of plausible spectroscopic-selection and augmentation realisations, and combining them through Dirichlet-weighted draws, we construct an experiment-marginalised ensemble of tomographic redshift distributions that encodes both statistical and systematic uncertainty. This ensemble naturally yields priors on commonly used nuisance descriptions, such as per-bin mean shifts and width rescalings, together with their cross-bin covariance. However, our results also show that these low-dimensional parametrisations are not generically sufficient. In particular, for source bins with extended or asymmetric tails, experiment-to-experiment variation cannot be captured fully by shifting the mean and affinely rescaling the width alone. This motivates the release of bias-corrected products not only in shift-only and shift-plus-scale form, but also as full-shape ensemble products, so that future cosmological forecasts can test directly how much information is lost when redshift uncertainty is compressed to a small number of nuisance parameters.

The present study remains a controlled forecasting analysis, and its dominant limitations are now clear. The largest residual systematic arises from mismatch between the simulated augmentation population and the target photometric population, especially in the colour--redshift relation of faint, high-redshift galaxies. Additional limitations arise from the stylised treatment of spectroscopic selection, the absence of imaging-level systematics such as blending, variable depth, and PSF-driven measurement biases, and the known shortcomings of the current estimators: posterior stacking does not satisfy the consistency conditions identified by \citet{2021PhRvD.103h3502M}, and the \texttt{Hybrid} combination remains ad hoc. While SOM-based augmentation partially mitigates these issues by improving feature-space completeness and regularising the effective training prior, none of the estimators carries formal optimality guarantees. These effects must be addressed in future work. A natural next step is therefore to deploy this framework on early LSST data using overlap spectroscopy, Deep Drilling Fields, and external deep imaging, while extending the present single-SOM implementation towards a multi-SOM architecture that can transfer information across surveys with different depths and filter coverage. In parallel, broader simulation families and imaging-level forward modelling will be required to turn augmentation from a single-simulation correction into a principled uncertainty model for realistic LSST calibration.

Overall, this work shows that SOM-based data augmentation, when combined with hybrid conditional-density estimation and experiment-level marginalisation, provides a practical and extensible route to improved calibration of tomographic redshift distributions. The main lesson is twofold: augmentation can substantially reduce calibration bias and stabilise inference against incomplete spectroscopy, but robust Stage-IV redshift calibration also requires uncertainty products that retain correlated, shape-level freedom beyond standard nuisance models. This framework therefore provides a concrete set of shift-only, shift-plus-scale, and full-shape ensemble products for forecasting studies, together with a methodological foundation for future LSST redshift calibration.